\section{Introduction}
%should I say high-fidelity or good approximations to
Preparing high-fidelity ground states is a central task in quantum computing, with applications in quantum chemistry \cite{kandala2017hardware}, quantum optimization \cite{farhi2014quantum, moll2018quantum}, and quantum machine learning \cite{farhi2018classification, cong2019quantum, peruzzo2014variational, cerezo2021variational, koczor2022quantum, koczor2022quantum2}. Furthermore, the efficiency of many fault-tolerant algorithms (e.g. phase estimation \cite{wang2022state}) hinges on the quality of their input state.

The Variational Quantum Eigensolver (VQE), a prominent example of a variational quantum algorithm (VQA), is a well-studied, hybrid, quantum-classical algorithm proposed to prepare the ground state of a Hamiltonian $H$ by variationally optimizing the parameters of an ansatz circuit $U(\theta)$ \cite{cerezo2021variational, peruzzo2014variational}. Since the ansatz structure is typically fixed and only the parameters $\theta$ are varied during training, choosing the right ansatz structure is crucial as it determines the family of expressible states \cite{kandala2017hardware, romero2018strategies, wecker2015progress}. It has been observed in the literature that fixed ansatz structures lead to a fundamental trade-off between expressivity and trainability \cite{holmes2022connecting}: highly expressive ansätze tend to exhibit barren plateaus of exponentially small gradients \cite{mcclean2018barren, larocca2024review, cerezo2021cost, wang2021noise}, whereas simpler ansätze may lack the ability to approximate the target state accurately. Trainability issues are further exacerbated by the presence of poor local minima \cite{anschuetz2022quantum, bittel2021training}. %In practice, accurate energy (and, when applicable, gradient) estimation requires many repeated circuit evaluations, which substantially increases the measurement overhead. 
Consequently, generic applications of both gradient-based and gradient-free VQAs may demand a prohibitively large measurement budget, constituting a major bottleneck for scalability on realistic hardware. 

Resource minimization becomes even more critical in early fault-tolerant regimes, where deeper circuits may be feasible but logical measurement rates remain limited \cite{zimboras2025myths}. These considerations have motivated the development of alternative state-preparation strategies that retain flexibility, avoid expensive gradient estimation, and aim to extract maximal information from each quantum measurement \cite{boyd2022training, hwang2025preparing, puig2026warm,grimsley2019adaptive, nakanishi2020sequential}. 
Classical shadows \cite{huang2020predicting} are particularly relevant in this context, as they enable the simultaneous estimation of many expectation values from a modest number of measurements. Specifically, given a quantum state $\rho$, one performs a measurement in a randomly chosen basis (e.g., random local Pauli or Clifford measurements) and constructs from each outcome a classical snapshot (a shadow) $\hat{\rho}$. These snapshots can be used to estimate expectation values of a wide range of observables. Remarkably, for fixed accuracy and failure probability, the number of required snapshots scales only logarithmically in the number $M$ of target observables, while generally growing exponentially with their locality under random local Pauli measurements. 

In this work, we introduce a resource-efficient, greedy, adaptive ground-state preparation algorithm that we term \emph{Shadow Enhanced Greedy Quantum Eigensolver} (SEGQE). Given an input state $\ket{\psi_0}$ and a Hamiltonian $H$, SEGQE iteratively constructs a circuit starting from the identity. At each step, a batch of classical shadows of the current state is collected and used to estimate the energy decrease induced by a set of candidate gate additions. The gate yielding the largest predicted energy decrease is appended to the circuit. Repeating this procedure progressively refines the circuit and yields a state that approximates the ground state of $H$. 

\begin{figure*}[htbp]
	\begin{center}
		\tikzfig{algo}
	\end{center}
	\caption{Schematic overview of the Shadow-Enhanced Greedy Quantum Eigensolver (SEGQE). At iteration $k$, a quantum computer prepares the state $\ket{\psi_k}=C_k\ket{\psi_0}$ and $N$ independent classical shadows are collected by applying random unitaries $V\in \mathcal{U}$ and measuring in the computational basis. In this work, $\mathcal{U}=\{I, H, S^\dagger H\}^{\otimes n}$, corresponding to uniformly random single-qubit Pauli measurements. The shadows are then used to estimate all Pauli expectation values needed to evaluate, for each candidate gate $U_j(\theta)\in\mathcal{G}$, the energy decrease $\Delta E_{k,j}(\theta)$. A classical optimizer finds $(j^*, \theta^*_{j^*})$ maximizing the expected decrease and appends $U_{j^*}(\theta^*_{j^*})$ to obtain $C_{k+1}$. The procedure repeats until convergence or until a predefined circuit depth is reached.}
	\label{fig:SEGQE}
\end{figure*}

The central advantage of SEGQE lies in its measurement efficiency. At each iteration, a single batch of classical shadows can be reused to evaluate a large set of candidate gate additions entirely in classical post-processing, yielding a per-iteration sample complexity that scales only logarithmically with the size of the candidate set. Furthermore, this design offloads the computationally expensive search over gate additions to classical hardware, enabling efficient parallelization and utilization of high-performance computing resources. While the circuit depth increases with successive gate additions, the per-iteration measurement cost grows only mildly with system size, and the information extracted per shot retains the asymptotic optimality guarantees of classical shadows. This feature is particularly well suited to early fault-tolerant regimes, where deeper circuits than typical VQA ansätze may be feasible, but relatively slow logical measurement rates demand shot-frugal techniques. In such settings, minimizing the measurement budget while exploiting classical computational power is crucial. SEGQE therefore provides a promising approach for applications such as initial state preparation for fault-tolerant phase estimation. 

A variety of iterative state-preparation methods have been proposed in the literature~\cite{motta2020determining, grimsley2019adaptive, feniou2025greedy} (see \Cref{appendix:related_works}). Here, we introduce a general classical-shadows-based framework that accommodates large sets of arbitrary local unitary gate additions and admits rigorous worst-case bounds on the per-iteration sample complexity with logarithmic dependence on the number of candidate gates. Although greedy strategies are not guaranteed to converge to the exact ground state in general, our numerical experiments demonstrate that SEGQE prepares high-fidelity ground-state approximations across a range of Hamiltonians.

The remainder of this article is structured as follows. In \cref{sec:procedure} we give a detailed description of the SEGQE algorithm, and in \cref{sec:guarantees} we provide rigorous worst-case upper bounds on its per-iteration sample complexity. In \cref{sec:numerics} we benchmark SEGQE on transverse-field Ising models and random local Hamiltonians, and study how the choice of gate set impacts convergence and final accuracy. Finally, in \cref{sec:conclusion} we reflect on our results and provide directions for future research. 

